% =====================================================================
% Appendice -- Glossario e acronimi
% =====================================================================
\chapter{Glossario e acronimi}
\label{app:glossario}

\begin{nota}[Fonte]
Rielaborato in italiano dal glossario distribuito con il primo laboratorio (\emph{Glossary -- acronyms,
laws \& regulations}), più alcune voci emerse a lezione. I riferimenti normativi completi sono nei
capitoli del manuale; qui una riga per voce.
\end{nota}

\section{Acronimi}
\begin{longtable}{>{\bfseries}p{0.17\textwidth}p{0.76\textwidth}}
\toprule
\normalfont\textbf{Sigla} & \textbf{Significato} \\
\midrule
\endhead
ACN & Agenzia per la Cybersicurezza Nazionale: autorità italiana competente per la NIS2 e autorità
  nazionale di certificazione; ospita lo CSIRT Italia. \\
BCR & \emph{Binding Corporate Rules}: regole vincolanti d'impresa per i trasferimenti infragruppo
  (art.~47 GDPR). \\
CB & \emph{Circumstances of the Breach}: fattore additivo della metodologia ENISA (perdita di
  riservatezza, integrità, disponibilità, intento malevolo). \\
CIA & \emph{Confidentiality, Integrity, Availability}: le tre proprietà di sicurezza; i tipi di
  violazione le rispecchiano. \\
CJEU & Corte di giustizia dell'UE (Lussemburgo): interpretazione vincolante del diritto UE. Da non
  confondere con la CEDU. \\
CRA & \emph{Cyber Resilience Act}, Reg.~(UE) 2024/2847. \\
CSA & \emph{Cybersecurity Act}, Reg.~(UE) 2019/881. \\
CSIRT & \emph{Computer Security Incident Response Team}: riceve le segnalazioni NIS2 e CRA (in Italia
  CSIRT Italia, presso l'ACN). \\
DORA & \emph{Digital Operational Resilience Act}, Reg.~(UE) 2022/2554 (settore finanziario). \\
DPA & \emph{Data Protection Authority} (autorità di controllo). Attenzione: anche \emph{Data
  Processing Agreement}, il contratto dell'art.~28 GDPR. \\
DPC & \emph{Data Processing Context}: punteggio base ENISA (da 1 a 4) che dipende dal tipo di dati. \\
DPIA & Valutazione d'impatto sulla protezione dei dati (art.~35 GDPR), obbligatoria per i trattamenti
  probabilmente ad alto rischio. \\
DPO & \emph{Data Protection Officer}, responsabile della protezione dei dati (artt.~37--39 GDPR). \\
ECtHR / CEDU & Corte europea dei diritti dell'uomo (Strasburgo, Consiglio d'Europa): non è
  un'istituzione UE. \\
EDPB & Comitato europeo per la protezione dei dati: riunisce le autorità di controllo e adotta linee
  guida (es. 9/2022 sulla notifica dei data breach). \\
EDPS & Garante europeo della protezione dei dati, competente per le istituzioni UE. \\
EI & \emph{Ease of Identification}: moltiplicatore ENISA (0,25; 0,5; 0,75; 1). \\
ENISA & Agenzia dell'UE per la cybersicurezza; mandato permanente con il CSA. \\
ESAs & Autorità europee di vigilanza (EBA, ESMA, EIOPA): sorveglianza ai sensi di DORA. \\
EUCC & Primo schema europeo di certificazione adottato con il CSA, basato sui Common Criteria. \\
FRIA & Valutazione d'impatto sui diritti fondamentali (AI Act), per alcuni \emph{deployer} di sistemi
  ad alto rischio. \\
GPAI & Modelli di IA per finalità generali (AI Act). \\
ISMS & Sistema di gestione della sicurezza delle informazioni (ISO/IEC~27001). \\
LED & Direttiva (UE) 2016/680 sul trattamento dei dati in ambito penale e di polizia: fuori dal GDPR. \\
PDCA & \emph{Plan--Do--Check--Act}: ciclo di miglioramento continuo dei sistemi di gestione ISO. \\
PDE & \emph{Product with Digital Elements}: prodotto con elementi digitali, ambito del CRA. \\
RTO / RPO & \emph{Recovery Time / Recovery Point Objective}: obiettivi di ripristino e di perdita
  massima di dati. \\
SBOM & \emph{Software Bill of Materials}: distinta dei componenti software (CRA). \\
SCCs & Clausole contrattuali tipo approvate dalla Commissione per i trasferimenti (art.~46 GDPR). \\
SE & \emph{Severity}: risultato ENISA, $SE = DPC \times EI + CB$. \\
SoA & \emph{Statement of Applicability}: elenco dei controlli dell'Annex~A ISO/IEC~27001 applicati
  o esclusi, con giustificazione. \\
TIA & \emph{Transfer Impact Assessment}: valutazione richiesta con SCC e BCR dopo Schrems~II. \\
TOMs & Misure tecniche e organizzative (artt.~24, 25, 32 GDPR). \\
\bottomrule
\end{longtable}

\section{Norme e strumenti}
\begin{longtable}{>{\bfseries}p{0.27\textwidth}p{0.66\textwidth}}
\toprule
\normalfont\textbf{Strumento} & \textbf{In una riga} \\
\midrule
\endhead
Carta dei diritti fondamentali UE, artt.~7--8 & Rispetto della vita privata e protezione dei dati
  personali come diritto fondamentale autonomo. \\
Direttiva 95/46/CE & Direttiva madre sulla protezione dei dati, abrogata dal GDPR. \\
Direttiva ePrivacy 2002/58/CE & Riservatezza nelle comunicazioni elettroniche; primo obbligo di
  notifica dei breach (operatori telecom); cookie. \\
GDPR, Reg.~(UE) 2016/679 & Protezione dei dati personali; applicabile dal 25~maggio~2018. \\
NIS2, Dir.~(UE) 2022/2555 & Soggetti essenziali e importanti, misure di gestione del rischio,
  notifica degli incidenti significativi. \\
DORA, Reg.~(UE) 2022/2554 & Resilienza operativa digitale del settore finanziario; \emph{lex
  specialis} rispetto alla NIS2. \\
CSA, Reg.~(UE) 2019/881 & Mandato permanente di ENISA; quadro europeo di certificazione. \\
CRA, Reg.~(UE) 2024/2847 & Requisiti di cybersicurezza dei prodotti con elementi digitali;
  marcatura CE. \\
AI Act, Reg.~(UE) 2024/1689 & Regolamento sull'IA basato sul rischio (pratiche vietate, alto
  rischio, trasparenza, GPAI). \\
Convenzione 108 & Primo trattato internazionale vincolante sulla protezione dei dati (Consiglio
  d'Europa, 1981). \\
Convenzione di Budapest & Convenzione sul cybercrime (2001): reati informatici armonizzati e
  cooperazione. \\
D.lgs.~196/2003 & Codice privacy italiano, adeguato al GDPR dal D.lgs.~101/2018. \\
D.lgs.~138/2024 & Recepimento italiano della NIS2; autorità competente l'ACN. \\
ISO/IEC~27001:2022 & Requisiti dell'ISMS; certificabile; Annex~A con 93 controlli. \\
ISO/IEC~27002 & Linee guida di attuazione dei controlli dell'Annex~A. \\
ISO/IEC~27701 & Estensione privacy dell'ISMS (PIMS). \\
\bottomrule
\end{longtable}
